\chapter{Introduction}

\section{Motivation}

RISC-V \cite{RISC-VInternational2025} is a modern and extensible instruction set architecture often used in systems-on-chip,
especially, FPGA-based ones \cite{LiteXGitHub}, such that it is both extensible and energy-efficient.
RISC-V's instruction set architecture is designed with custom extensions in mind,
so designing custom extensions would be a straightforward process.
It is, however, a difficult task to integrate a custom extension into the RISC-V specification,
and even more so to implement an extension for the multitude of RISC-V processor implementations.
With such issues, and the fact that most FPGA-bases SOCs offer plenty
of space to also implement an extension,
in mind, Jan Gray et al. designed the custom composable extensions (CFU) \cite{Gray2022CFU} extension.
It allows for an implementation of a specific extension to be used across multiple RISC-V CPUs,
so as to use the single extension implementation across a multitude of different RISC-V CPUs.
CFU allows for a single extension to be connected to a hardware thread.
It also allows passing a "function index" to an extension, such that up to 255 different functions may be implemented.

There is also a playground for desigining the CFU extensions on a ready-made SOC called
the CFU playground. It allows the users to implement a
CFU and write a main script that runs and invokes the implemented extension.
CFU is a useful interface for implementing a single extension,
but when combining extensions, especially those,
that may have multiple different functions and utilize the function id, it becomes incredibly
difficult and cumbersome to switch between them.
Because of that, Gray et al. also designed a newer version of composable
custom extensions - the CXU \cite{Gray2025CXU}.
This implementation allows for switching between multiple custom extensions,
but, as of now, there is not yet an implementation of it
on any of the RISC-V implementations, nor an environment for developing custom extensions.

\section{Contributions}

The goal of this work is to implement the CXU RISC-V extension for the Vexii
RISC-V Processor and then to set up a development environment, or sandbox,
to simplify creating projects that utilize custom CXU compute units.

The contribution of this work is the sandbox that gives the user the ability
to implement custom extensions, as well as test and benchmark them.

To validate the effectiveness of the CXU sandbox, the research includes
the implementation and benchmarking of practical example extensions such as AES and GZIP.
The AES extension serves as a primary case study for accelerating cryptographic workloads,
demonstrating how modular hardware can achieve substantial runtime improvements
compared to software-only reference algorithms.

These examples collectively serve as proof of concept for the sandbox’s ability to
host diverse compute units. The results show that even complex algorithms
can be integrated and benchmarked efficiently, providing measurable gains
in instruction efficiency and overall cycle counts within a unified development environment.


\section{Structure Of The Thesis}

\begin{itemize}
    \item Chapter 2 surveys existing implementations of similar custom extensions sandboxes and custom instruction interfaces, as well as works that utilize those extensions to achieve some compute advantages.
    \item Chapter 3 provides the theoretical background for how custom SoC frameworks, the Vexii RISC-V and CXU work.
    \item Chapter 4 formulates the goals of this work in more detail and describes the methodology used in the thesis. This chapter covers how the CXU extension is implemented for Vexii RISC-V, how it's integrated into litex, the sandbox environment, example extensions and testing / benchmarking methodology used for this thesis.
    \item Chapter 5 covers the results of the implementation of the sandbox and the example extensions, tests for the CXU interface base functionality and benchmarking results.
    \item Chapter 6 summarizes the results achieved in this work and discusses potential improvements for the sandbox.
\end{itemize}
